\begin{abstract}
Electricity-market data are graded and released field by field, yet the inference risk of a field depends on what else an
adversary holds: fields that reveal little alone can jointly expose unit outputs, private offers, line flows or load. We
define field-level risk from set-level inferability, the out-of-sample accuracy of attackers retrained on a set of fields.
The budgeted marginal risk $M^{(K)}$ of a field is its largest inferability gain over any background of at most $K$ fields,
relative to data already public; its excess over single-field risk is the background amplification $\Gamma^{(K)}$.
$M^{(K)}$ is a budget-constrained marginal contribution index and the minimal threshold-free safety margin, and it flags
exactly the members of minimal unsafe combinations of at most $K+1$ fields. To evaluate the many field sets this requires,
we propose an amortized attacker---masked pre-training with column reconstruction, features of an untrained network, and a
closed-form readout per field set---whose scan is certified by retraining a few backgrounds per field. In a controlled
market case targeting a unit's private offer markup, every price field has zero risk alone and all of its risk arises over
the unit's output, ranked and congestion-dependent as dispatch physics implies. On four power data sets with rule-based
field roles and $1.2\times10^5$ exhaustively retrained field sets, a typical field's average contextual gain is a third of
its worst case, and single-field grading recovers 13--17\% of the critical fields, whereas scan--certify recovers up to 91\%
without false positives. The amortized attacker outperforms linearized, fine-tuned and surrogate estimators and certifies as
much of the risk as a tabular foundation model at one to two orders of magnitude lower cost per field set.
\end{abstract}
